% sections/03_system_design.tex

\section{System Design}
\label{sec:system}

\subsection{Architecture}
\label{sec:arch}

Figure~\ref{fig:architecture} shows the three-layer architecture of the
Agentic Chemometrician.
The top layer is \texttt{chemometrics\_mcp}, an MCP server implemented in
Python.
It imports from \texttt{mcp.server} and \texttt{mcp.types}, dispatches JSON
tool calls, and serializes all responses as \texttt{TextContent} JSON.
Thin \emph{tool wrappers} in \texttt{src/chemometrics\_mcp/tools/} translate
incoming JSON dictionaries into typed request contracts and pass them to the
core layer.

The middle layer is \texttt{chemometrics\_mcp.core}, which contains all
deterministic analysis logic: dataset loading, preprocessing, modeling,
validation, planning, fallback recommendation, method memory, and report
generation.
\textbf{No file in \texttt{core/} imports from \texttt{mcp}}.
This strict separation means the core can be tested, audited, and replayed
with direct Python imports---no MCP transport is required---and its outputs
are guaranteed to be identical across runs at the same seed and input.

The bottom layer is \texttt{chemometrics\_contracts}, a standalone Python
package of frozen dataclasses that defines the typed schema for every value
that crosses the tool boundary.
Both the MCP layer and the core layer import from contracts; neither direction
creates a circular dependency.

\begin{figure}[ht]
  \centering
  % TODO: replace with actual diagram PDF
  \IfFileExists{figures/fig1_architecture.pdf}{%
    \includegraphics[width=0.85\linewidth]{figures/fig1_architecture}%
  }{%
    \fbox{\parbox{0.9\linewidth}{\centering\vspace{2cm}\TODO{figure: fig1\_architecture}\vspace{2cm}}}%
  }
  \caption{
    System architecture of the Agentic Chemometrician.
    The MCP server layer (top) exposes JSON tool endpoints;
    the deterministic core layer (middle) contains all computation;
    the contracts layer (bottom) defines the typed dataclass schema shared
    by both.
    HITL approval gates (hatched boxes) appear at plan proposal,
    fallback selection, and final model selection under validation warnings.
    \textbf{[TODO: replace stub with actual diagram]}
  }
  \label{fig:architecture}
\end{figure}

\subsection{Pipeline}
\label{sec:pipeline}

The end-to-end analysis pipeline follows a fixed seven-stage sequence:
\textbf{inspect} $\to$ \textbf{propose} $\to$ \textbf{run} $\to$
\textbf{validate} $\to$ \textbf{select} $\to$ \textbf{interpret} $\to$
\textbf{report}.

\begin{enumerate}
  \item \textbf{Inspect} (\texttt{inspect\_dataset}).
    Load a spectral file (Excel or plain-text (.txt) format), infer or override modality,
    detect candidate label and group columns, enforce value-domain constraints,
    and return a \texttt{DatasetInspection} summary with axis statistics and
    data-quality warnings.

  \item \textbf{Propose} (\texttt{propose\_analysis\_plan}).
    Convert a \texttt{DatasetInspection} into a bounded \texttt{AnalysisPlan}:
    task type, preprocessing candidates appropriate to the modality,
    validation strategy (with GroupKFold or LOGO recommended when candidate
    group columns are detected), and initial model families.
    The plan is returned with a human-readable summary that the operator
    must approve before execution proceeds.
    \textbf{HITL gate: plan approval is required.}

  \item \textbf{Run} (\texttt{run\_analysis}).
    Execute the approved \texttt{AnalysisPlan} against a \texttt{SpectralDataset}.
    Only tasks present in the approved plan are executed.
    All cross-validation, preprocessing, model fitting, and metric computation
    happens in the core layer with a fixed random seed.
    Outputs are saved as JSON and PNG artifacts in a timestamped run directory.

  \item \textbf{Validate} (\texttt{validate\_results}).
    Run scientific reliability checks on the \texttt{AnalysisRun} output:
    replicate leakage, group leakage risk, class imbalance, small-sample
    warnings, split instability, suspicious metrics, modality--preprocessing
    consistency, and regression target leakage.
    Returns a \texttt{ValidationSummary} with severity-labeled warnings and
    a boolean \texttt{passed} flag.

  \item \textbf{Select} (\texttt{select\_best\_model}).
    Compare candidate models on performance, validation reliability,
    interpretability, stability, complexity, and task suitability.
    Distinguishes the \emph{best-measured} model from the
    \emph{best-defensible} model when validation warnings exist.
    \textbf{HITL gate: approval required when validation warnings are present.}

  \item \textbf{Interpret} (\texttt{interpret\_results}).
    Summarize feature or wavelength importance from tool-produced outputs.
    Compares evidence across models and separates model evidence from chemical
    conclusions.
    Flags unstable or weak interpretations via explicit warning codes.

  \item \textbf{Report} (\texttt{generate\_report}).
    Produce a human-readable and agent-readable final report from run
    artifacts, including metrics, figures, validation warnings, caveats,
    interpretations, next-step recommendations, and a human-review checklist.
\end{enumerate}

\subsection{Typed Contracts}
\label{sec:contracts}

All values that cross the tool boundary are instances of
\texttt{SerializableContract}, a base class that wraps Python's
\texttt{dataclasses.asdict} for lossless JSON round-trips.
Every contract dataclass is frozen (\texttt{frozen=True}) and uses
\texttt{slots=True} where available, preventing accidental mutation.

The contract library contains 30 frozen dataclasses over a shared
SerializableContract base: 18 domain contracts
(\texttt{SpectralDataset}, \texttt{AnalysisResult}, \texttt{AnalysisPlan},
\texttt{ValidationSummary}, \texttt{MethodMemory}, etc.) and 12
request/response wrappers (11 \texttt{*Request} classes plus
\texttt{ToolResponse[T]}).

The generic \texttt{ToolResponse[T]} wrapper unifies all tool outputs:
it carries \texttt{ok: bool}, a typed \texttt{payload}, a list of
\texttt{ValidationWarning} objects, artifact references, run metadata, and an
optional error message.
Every error path in the system returns a \texttt{ToolResponse} with
\texttt{ok=False} rather than raising an exception to the MCP transport layer.

\texttt{ValidationWarning} is itself a contract, with fields \texttt{code}
(a short machine-readable identifier), \texttt{message} (human-readable
explanation), \texttt{category}, \texttt{severity}, and
\texttt{affected\_stage}.
Warning codes are the primary API for downstream tools: \texttt{select\_best\_model}
reads \texttt{requires\_human\_approval} from \texttt{ModelSelectionRecommendation}
based on whether \texttt{ValidationWarning} codes include
\texttt{replicate\_leakage} or \texttt{group\_leakage\_risk}.

\subsection{Declarative Modality Registry}
\label{sec:registry}

Prior versions of the codebase contained scattered \texttt{if modality == "FTIR"}
branches in the preprocessing, validation, planning, and reporting modules.
We replaced these with a \emph{declarative modality registry}: a dictionary of
\texttt{ModalityProfile} frozen dataclasses, one per supported modality.

Each \texttt{ModalityProfile} records:
the canonical modality name and aliases;
the axis unit (\texttt{"cm-1"}, \texttt{"nm"});
the conventional axis direction (ascending/descending);
an axis inference range used to detect the modality from data (e.g.\ 700--2500~nm
for NIR, 400--4000~cm\textsuperscript{-1} for FTIR);
the data representation (\texttt{"gridded"} for continuous spectra,
\texttt{"peaklist"} for sparse MS data);
the valid intensity range and out-of-range policy (\texttt{"clip"} for FTIR
percent-transmittance, \texttt{"allow"} for NIR absorbance);
and the preprocessing candidates appropriate to that modality.

The registry currently supports FTIR, NIR, and Raman.
Adding a new gridded modality requires a single \texttt{ModalityProfile} entry
in the registry dictionary; no edits to any downstream module are required.
Sparse modalities (MS, NMR) require a peaklist resampler not yet implemented;
the \texttt{representation} field is present to make this distinction explicit.

The \texttt{infer\_modality(axis)} function resolves a spectrum's axis array
to a modality by finding the narrowest enclosing axis inference range,
preferring more specific matches (NIR's 700--2500~nm span is preferred over
FTIR's 400--4000~cm\textsuperscript{-1} span when the axis falls in the NIR range).
Raman shift, which overlaps the FTIR wavenumber range, cannot be inferred from
the axis alone and must be declared explicitly.

\subsection{Human-in-the-Loop Gates}
\label{sec:hitl}

Three HITL gates are hard-coded in the pipeline:

\begin{enumerate}
  \item \textbf{Plan approval.} The operator must approve the
    \texttt{AnalysisPlan} returned by \texttt{propose\_analysis\_plan} before
    \texttt{run\_analysis} is called.
    The plan includes a \texttt{human\_readable\_plan} string summarizing the
    proposed task, preprocessing, validation strategy, and model families.

  \item \textbf{Fallback acceptance.} When a model fails, \texttt{recommend\_next\_model}
    always sets \texttt{requires\_human\_approval=True} on the
    \texttt{NextModelRecommendation} it returns.
    The rationale text explicitly states that ``comparability may be limited ---
    human review required before accepting fallback results.''
    The gate cannot be bypassed programmatically.

  \item \textbf{Final model selection under warnings.} \texttt{select\_best\_model}
    sets \texttt{requires\_human\_approval=True} on the returned
    \texttt{ModelSelectionRecommendation} whenever the \texttt{ValidationSummary}
    contains warnings.
    The field distinguishes the best-measured model (highest metric) from the
    best-defensible model (highest metric conditional on passing all validation
    checks), and requires a human to choose between them when they differ.
\end{enumerate}
